% Slide — Clássicas: mantidas intactas
\begin{frame}{As três clássicas originais: nada mudou}
  \centering
  \tabstyle\footnotesize
  \begin{tabular}{@{}lll@{}}
    \toprule
    \textbf{Política} & \textbf{Regra} & \textbf{Parâmetros} \\
    \midrule
    EOQ & $Q^\ast=\sqrt{2D K/h}$; pede quando $I{+}\text{pipeline}\leq \mathrm{ROP}$ &
      $\mathrm{ROP}=\mu L + z\sigma\sqrt{L}$ \\
    $(s,S)$ & Banda min-max: pede até $S$ quando posição $\leq s$ &
      $s=\mu L{+}z\sigma\sqrt{L}$; $S=s{+}Q_{\mathrm{EOQ}}$ \\
    Jornaleiro & $Q^\ast=\hat\mu+z\hat\sigma$ (quantil crítico fixo) &
      $z=1{,}28$ (SL alvo 90\%) \\
    \bottomrule
  \end{tabular}
  \vspace{1em}

  \begin{block}{Por que preservar}
    \small São o \highlight{baseline do varejo} e o denominador de todas as
    comparações já reportadas no Capítulo 5. Mudar sua fórmula invalidaria
    a leitura de qualquer resultado anterior que as usa como referência.
  \end{block}

  \note{
    Estas três são as únicas políticas do portfólio original de 12 que
    permanecem byte-a-byte equivalentes em lógica -- só a chamada ao
    simulador por baixo delas mudou (do InventoryEnv escalar, que
    continua existindo e sendo usado por elas, sem alteração). A decisão
    de mantê-las intocadas foi deliberada: são o piso de comparação usado
    em todo o Capítulo 5 da proposta, e qualquer mudança nelas quebraria a
    comparabilidade com os números já reportados.
  }
\end{frame}
